% Part: sets-functions-relations
% Chapter: relations
% Section: orders

\documentclass[../../../include/open-logic-section]{subfiles}

\begin{document}

\olfileid{sfr}{rel}{ord}
\olsection{ترتيبونه}

\begin{explain}
په ډېرو پرتلو کښې ځينې څيزونه، له يوه پلوه، تر نورو ``کوچني''،
ورسره ``برابر'' يا ترې ``لوي'' بولو۔ په دې کښې د \emph{ترتيب}
اړيکې شاملې دي۔ خو د ترتيبي اړيکو بېل ډولونه شته۔ مثلاً، ځينې
لازموي چې هر دوه څيزونه د پرتلې وړ وي، خو نورې داسې نۀ دي۔ ځينې
عينيت هم رانغاړي (لکه $\le$)، او ځينې ئې باسي (لکه $<$)۔ دلته د
ډولونو وېش لرل راسره مرسته کوي۔
\end{explain}

\begin{defn}[مقدماتي ترتيب]
هغه اړيکه چې انعکاسي او انتقالي دواړه وي، \emph{مقدماتي ترتيب}
بللے شي۔
\end{defn}

\begin{defn}[جزوي ترتيب]
هغه مقدماتي ترتيب چې ضد متناظر هم وي، \emph{جزوي ترتيب} بللے شي۔
\end{defn}

\begin{defn}[خطي ترتيب]\ollabel{def:linearorder}
هغه جزوي ترتيب چې تړلے هم وي، \emph{کلي ترتيب} يا \emph{خطي ترتيب}
بللے شي۔
\end{defn}

\begin{ex}
هر خطي ترتيب جزوي ترتيب هم دے، او هر جزوي ترتيب مقدماتي ترتيب
هم دے؛ خو برعکس خبرې سمې نۀ دي۔ په $A$ باندې ټولشموله اړيکه
مقدماتي ترتيب دے، ځکه انعکاسي او انتقالي ده۔ خو که $A$ له يوه
!!{element} څخه زيات ولري، نو ټولشموله اړيکه ضد متناظره نۀ ده،
او له دې امله جزوي ترتيب هم نۀ دے۔
\end{ex}

\begin{ex}
په $\Bin^*$ باندې د \emph{تر بل نۀ اوږدېدو} اړيکه
$\preccurlyeq$ وګورئ: $x \preccurlyeq y$ هله او يوازې هله چې
$\len{x} \le \len{y}$ وي۔ دا مقدماتي ترتيب دے (انعکاسي او انتقالي
دے)، او تړلے هم دے؛ خو جزوي ترتيب نۀ دے، ځکه ضد متناظر نۀ دے۔
مثلاً، $01 \preccurlyeq 10$ او $10 \preccurlyeq 01$، خو
$01 \neq 10$۔
\end{ex}

\begin{ex}
د سټونو په يوه سټ باندې د $\subseteq$ اړيکه يو مهم جزوي ترتيب دے۔
دا په عمومي توګه خطي ترتيب نۀ دے، ځکه که $a \neq b$ وي او
$\Pow{\{a, b\}} = \{\emptyset, \{a\}, \{b\}, \{a,b\}\}$
وګورو، نو وينو چې $\{a\} \nsubseteq \{b\}$،
$\{a\} \neq \{b\}$ او $\{b\} \nsubseteq \{a\}$۔
\end{ex}

\begin{ex}
د \emph{بې له باقي وېشل کېدو} اړيکه داسې جزوي ترتيب راکوي چې خطي
ترتيب نۀ دے۔ د صحيح عددونو $n$ او $m$ دپاره، $n \mid m$ په دې
معنا ليکو چې $n$، $m$ (بې له باقي) وېشي؛ يعنې هله او يوازې هله
چې داسې صحيح عدد $k$ شته چې $m = kn$ وي۔ په $\Nat$ باندې دا جزوي
ترتيب دے، خو خطي ترتيب نۀ دے: مثلاً، $2 \nmid 3$ او $3 \nmid 2$
هم۔ که وېشمنتيا په $\Int$ باندې اړيکه وګڼو، يوازې مقدماتي ترتيب
دے، ځکه ضد متناظره نۀ ده: $1 \mid -1$ او $-1 \mid 1$، خو
$1 \neq -1$۔
\end{ex}

\begin{ex}
د لړيو په سټ $A^*$ باندې د \emph{غځونې} اړيکه دا ده:
$s \sqsubseteq s'$ هله او يوازې هله چې $s = \emptyseq$ (تشه لړۍ)
وي، $s = s'$ وي، يا $s = \tuple{s_1, \dots, s_n}$ او
$s' = \tuple{s_1, \dots, s_n, s_{n+1}, \dots, s_m}$ وي۔ که
$s \sqsubseteq s'$ وي، دا هم وايو چې $s$ د $s'$ \emph{پيلنۍ برخه}
ده۔ په $A^*$ باندې د غځونې اړيکه جزوي ترتيب دے، خو خطي ترتيب نۀ
دے؛ مثلاً، که $a \neq b$ وي، نو $ab \not\sqsubseteq ba$ او
$ba \not\sqsubseteq ab$۔
\end{ex}

\begin{defn}[سخت ترتيب]
\emph{سخت ترتيب} هغه اړيکه ده چې نفي انعکاسي، نامتناظره او انتقالي وي۔
\end{defn}

\begin{defn}[سخت خطي ترتيب]\ollabel{def:strictlinearorder}
هغه سخت ترتيب چې تړلے هم وي، \emph{سخت کلي ترتيب} يا
\emph{سخت خطي ترتيب} بللے شي۔
\end{defn}

\begin{ex}
$\le$ هغه خطي ترتيب دے چې له سخت خطي ترتيب $<$ سره سمون خوري۔
$\subseteq$ هغه جزوي ترتيب دے چې له سخت ترتيب $\subsetneq$ سره
سمون خوري۔
\end{ex}

په $A$ باندې هر سخت ترتيب $R$ د قطر $\Id{A}$ په ورزياتولو، يعنې
د ټولو جوړو $\tuple{x, x}$ په ورزياتولو، جزوي ترتيب ګرځولے شو۔
(دې ته د $R$ \emph{انعکاسي بندښت} وايي۔) برعکس، له جزوي ترتيبه
په پيل کولو، د $\Id{A}$ په لرې کولو سخت ترتيب ترلاسه کولے شو۔
لاندې دوه نتيجې دا خبره دقيقه کوي۔

\begin{prop}\ollabel{prop:stricttopartial}
که $R$ په $A$ باندې سخت ترتيب وي، نو $R^+ = R \cup \Id{A}$
جزوي ترتيب دے۔ سربېره پر دې، که $R$ سخت خطي ترتيب وي، نو $R^+$
خطي ترتيب دے۔
\end{prop}

\begin{proof}
فرض کړئ $R$ سخت ترتيب دے؛ يعنې $R \subseteq A^2$ او $R$ نفي
انعکاسي، نامتناظر او انتقالي دے۔ $R^+ = R \cup \Id{A}$ کېږدئ۔
بايد وښيو چې $R^+$ انعکاسي، ضد متناظر او انتقالي دے۔

ښکاره ده چې $R^+$ انعکاسي دے، ځکه د هر $x \in A$ دپاره
$\tuple{x, x} \in \Id{A} \subseteq R^+$۔

د دې د ښودلو دپاره چې $R^+$ ضد متناظر دے، د تناقض له لارې د ثبوت
دپاره فرض کړئ $R^+xy$ او $R^+yx$، خو $x \neq y$۔ ځکه چې
$\tuple{x,y} \in R \cup \Id{A}$، خو
$\tuple{x, y} \notin \Id{A}$، نو خامخا $\tuple{x, y} \in R$
لرو، يعنې $Rxy$۔ همدارنګه، $Ryx$۔ خو دا له دې فرضيې سره تناقض لري
چې $R$ نامتناظر دے۔

د انتقالي والي د ثابتولو دپاره، فرض کړئ $R^+xy$ او $R^+yz$۔
که $\tuple{x, y} \in R$ او $\tuple{y,z} \in R$ دواړه وي، نو
$\tuple{x, z} \in R$، ځکه $R$ انتقالي دے۔ که داسې نۀ وي، نو يا
$\tuple{x, y} \in \Id{A}$، يعنې $x = y$، يا
$\tuple{y, z} \in \Id{A}$، يعنې $y = z$۔ په لومړي حالت کښې د
فرضيې له مخې $R^+yz$ لرو، او $x = y$، نو $R^+xz$۔ په دويم
حالت کښې هم همداسې ده۔ په هر حالت کښې $R^+xz$، نو $R^+$ انتقالي
هم دے۔

د ``سربېره پر دې'' برخې په اړه، فرض کړئ $R$ تړلے هم دے۔ نو د ټولو
$x \neq y$ دپاره يا $Rxy$ يا $Ryx$؛ يعنې يا
$\tuple{x, y} \in R$ يا $\tuple{y, x} \in R$۔ ځکه چې
$R \subseteq R^+$، دا خبره د $R^+$ دپاره هم سمه پاتې کېږي۔ نو
$R^+$ تړلے هم دے۔
\end{proof}

\begin{prop}\ollabel{prop:partialtostrict}
که $R$ په $A$ باندې جزوي ترتيب وي، نو
$R^- = R \setminus \Id{A}$ سخت ترتيب دے۔ سربېره پر دې، که $R$
خطي ترتيب وي، نو $R^-$ سخت خطي ترتيب دے۔
\end{prop}

\begin{proof}
دا د تمرين په توګه پرېښودل شوے دے۔
\end{proof}

\begin{prob}
د \olref[sfr][rel][ord]{prop:partialtostrict} ثبوت ورکړئ۔
\end{prob}

لاندې ساده نتيجه ثابتوي چې سخت خطي ترتيبونه د غړو له مخې د
برابرۍ په شان يو خاصيت پوره کوي:

\begin{prop}\ollabel{prop:extensionality-strictlinearorders}
که $<$ په $A$ باندې سخت خطي ترتيب وي، نو:
\[
  (\forall a, b \in A)((\forall x \in A)(x < a \liff x < b) \lif a = b).
\]
\end{prop}

\begin{proof}
فرض کړئ $(\forall x \in A)(x < a \liff x < b)$۔ که $a < b$ وي،
نو $a < a$، چې د $<$ له نفي انعکاسي والي سره تناقض لري؛ نو
$a \nless b$۔ بالکل په همدې ډول $b \nless a$۔ نو $a = b$، ځکه
$<$ تړلے دے۔
\end{proof}

\end{document}
